\chapter{Formulação e metodologia}\label{formulauxe7uxe3o-e-implementauxe7uxe3o}

\section{Modelo de propagação}\label{modelo-de-propagauxe7uxe3o}

Adotar explicitamente

\[\omega^2=c^2k^2+\left(\frac{m_gc^2}{\hbar}\right)^2,\qquad
f_g=\frac{m_gc^2}{h_{\mathrm P}},\qquad
\beta(f)=\frac{ck}{\omega}=\sqrt{1-\left(\frac{f_g}{f}\right)^2}.\]

Aqui \(h_{\mathrm P}\) é a constante de Planck. A velocidade de grupo é
\(v_g=c\beta\); a velocidade de fase é \(v_{\mathrm{ph}}=c/\beta\). Essa
distinção precisa ser preservada no cálculo dos fatores de resposta.
Para \(f<f_g\), o modo livre não é uma onda viajante com número de onda
real.

A etapa inicial trabalha no regime linear sobre fundo plano, apropriado
para estudar a resposta local e a propagação na linha de visada
Terra--pulsar. A geração cosmológica e a distribuição das amplitudes
serão hipóteses especificadas, não previsões completas de dRGT.
Fierz--Pauli será utilizado como referência linear de um campo massivo
saudável. Isso não equivale a resolver a emissão por binárias em uma
teoria não linear nem o mecanismo de Vainshtein.

\section{Polarizações e
vínculos}\label{polarizauxe7uxf5es-e-vuxednculos}

Calcular diretamente \(E_{ij}=R_{0i0j}\) a partir de \(h_{\mu\nu}\).
Para Fierz--Pauli, impor os vínculos de transversalidade
quadridimensional e traço nulo. Para a helicidade zero, na convenção

\[p_b=R_{0x0x}+R_{0y0y},\qquad p_l=R_{0z0z},\]

o benchmark linear de onda plana resulta em

\[\frac{p_l}{p_b}=-\left(\frac{f_g}{f}\right)^2.\]

Essa relação será uma verificação da implementação. Ela não é a fórmula
de um escalar genérico em qualquer teoria. Tampouco implica que todas as
amplitudes adicionais sejam universalmente suprimidas por \(m_g^2/f^2\).
A intensidade excitada depende do modelo e da fonte.

\section{Resposta e correlações de
PTA}\label{resposta-e-correlauxe7uxf5es-de-pta}

Derivar a variação de frequência dos pulsos e integrá-la para obter os
resíduos de temporização. A matriz de marés local, isoladamente, não
constitui a resposta de uma PTA. Para cada polarização independente
\(A\), construir \(\mathcal R_a^A(f,\hat\Omega,L_a;m_g)\), incluindo os
termos da Terra e do pulsar, e então

\[\Gamma_{ab}^{A}(f)=\mathcal N_A\int d^2\hat\Omega\,
\mathcal R_a^{A}(f,\hat\Omega)\mathcal R_b^{A*}(f,\hat\Omega).\]

Fixar e documentar \(\mathcal N_A\), as convenções espectrais e o
tratamento das autocorrelações. Comparar integração angular direta com
um segundo método ou limite analítico. Usar as expressões de
\textcite{liang2021} e as correções aplicáveis de \textcite{cordes2025}
como benchmarks, após compatibilizar normalizações.

O limite de Hellings--Downs será exigido no setor tensorial apropriado.
Não se deve exigir que o modelo massivo completo com amplitudes
escalares arbitrárias se torne RG apenas tomando \(m_g\to0\); o limite
com fontes e o desacoplamento dos setores adicionais exigem tratamento
próprio.

\section{Comparação central}\label{comparauxe7uxe3o-central}

Construir, a partir das mesmas realizações simuladas:

\begin{itemize}
\tightlist
\item
  \textbf{Análise com frequência explícita:} modelo de covariância ou
  espectros cruzados dependentes de \(f\), usando uma mesma massa em
  todas as frequências.
\item
  \textbf{Análise comprimida consistente:} aplicar ao modelo os mesmos
  pesos, janela e compressão utilizados nos estimadores simulados.
\item
  \textbf{Análise com frequência de referência:} substituir a
  dependência relevante por uma avaliação em \(f_{\mathrm{ref}}\),
  reproduzindo uma aproximação controlada.
\end{itemize}

Esquematicamente, a previsão comprimida correta tem a forma

\[\langle\widehat C_{ab}\rangle=\int df\,W_{ab}(f)S_h(f)\Gamma_{ab}(f;m_g),\]

em que \(W_{ab}\) inclui o estimador e a observação. Em geral, isso não
coincide com retirar \(\Gamma_{ab}(f_{\mathrm{ref}};m_g)\) da integral.
Não existe uma frequência efetiva universal independente do espectro e
dos pesos. A extensão com várias polarizações deve somar espectros e
respostas dos setores independentes.

Essa comparação separa a perda causada pela compressão da perda causada
por aproximar incorretamente o modelo comprimido. Também permite avaliar
covariância completa e diagonal sem pressupor que ignorar correlações
seja sempre conservador.

Será útil medir o erro do modelo em unidades do ruído, por exemplo
\(\Delta\mathbf C^\mathsf T\Sigma^{-1}\Delta\mathbf C\), além de erros
relativos nas ORFs. Isso aproxima a discussão da consequência
observacional: uma discrepância matemática grande em um canal pouco
sensível pode ter efeito estatístico pequeno.

\section{Experimentos numéricos e
inferência}\label{experimentos-numuxe9ricos-e-inferuxeancia}

Começar com PTAs sintéticas de 10--30 pulsares, duração de 4,5 e 15 anos
e uma seleção controlada de geometrias, cadências e ruído. Esses
tamanhos são escolhas de planejamento e poderão ser ajustados após medir
o custo computacional. Comparar espectro em lei de potência com um
espectro de baixa frequência modificado; acrescentar sinais comuns
monopolar e dipolar como testes de robustez.

Explorar massas pela razão adimensional \(f_gT\) e pela posição do
limiar nos canais observados. Se um canal estiver abaixo de \(f_g\),
modelar a ausência da contribuição viajante daquele setor. Não excluir
automaticamente todo o modelo por causa de um canal sem sinal
confirmado. Uma mistura entre setores massivos e não massivos exige
suporte espectral separado.

Para massa positiva, comparar escolhas a priori uniformes em \(m_g\) e
em \(m_g^2\), e uma escolha logarítmica com corte inferior explicitado.
Incluir RG como hipótese separada: uma priori contínua em \(\log m_g\)
não contém o ponto \(m_g=0\).

As amplitudes dos setores tensorial e escalar serão parâmetros de
população especificados. As deformações transversal e longitudinal do
mesmo modo escalar deverão permanecer correlacionadas. Os modos
vetoriais e a análise integral de uma segunda PTA serão extensões
condicionais ao progresso.

\clearpage
\section{Dados e ferramentas}\label{dados-e-ferramentas}

Usar Python com álgebra simbólica e rotinas numéricas verificáveis. Para
uma eventual análise completa de resíduos, aproveitar o
ENTERPRISE, mantido pelo NANOGrav \autocite{enterprise}, em vez de reconstruir toda a infraestrutura de
temporização. Fixar versões e sementes aleatórias.

Os dados públicos do NANOGrav \autocite{nanogravdata} e os produtos documentados pelos artigos
servirão como referências. Antes da etapa observacional, verificar quais
produtos contêm informação em frequência e covariâncias. Correlações
angulares já comprimidas não permitem recuperar informação espectral
perdida.

O artigo principal poderá ser concluído com simulações e benchmarks. A
aplicação a dados públicos será realizada se houver produtos adequados e
tempo, declarando as limitações de qualquer compressão ou covariância
indisponível. Uma análise completa dos resíduos não será substituída por
um ajuste de pontos de uma figura.
